\originalsection{行波方向与介质中的传播}\label{tb:wave-directions}
\begin{example}
\textbf{原13.12.1、13.12.2: 指定一列波后求配对磁场.} 真空中$\ee=E_0\cos(kz-\omega t)\uv x$, $k>0,\omega=ck$. 求传播方向和磁场; 再验证正弦场满足波动方程.
\end{example}
\begin{solution}
相位$kz-\omega t$固定的位置满足$z=(\omega/k)t+\text{常数}$, 波向$+z$. 因而配对磁场$\bb=(E_0/c)\cos(kz-\omega t)\uv y$, $\ee\times\bb$指$+z$. 法拉第的$y$分量为$\partial_zE_x=-\partial_tB_y$: 两边分别为$-kE_0\sin$和$-\omega(E_0/c)\sin$, 一致. 相似地安培--麦克斯韦给$kB_0=\omega\mu_0\eps E_0$. 二阶导数$\partial_z^2E=-k^2E$, $\partial_t^2E=-\omega^2E$, 所以$\partial_z^2E-c^{-2}\partial_t^2E=0$; 磁场同样满足. 这个代入核对还解释了$\omega=ck$条件为何不可省.
\end{solution}

均匀、线性、无损且非色散介质中, 无自由源的场方程把$\eps,\mu_0$替换为$\varepsilon,\mu$. 同样消元得$v=1/\sqrt{\mu\varepsilon}$, 单列平面行波$B_0=E_0/v$, 波阻抗$Z_m=\sqrt{\mu/\varepsilon}$. 稳定静止界面上, 连续边界条件要在所有时刻成立, 因而入射、反射、透射的频率相同, 波数和波长可变: $\lambda_m=v/f$. 这是光程相位的物理来源. 真空能流式中的$\mu_0$在此换成$\mu$, 能密度为$\varepsilon E^2/2+B^2/(2\mu)$; 强色散、吸收材料不能直接用这组简单能量式.

\originalsection{光强、动量与球面扩散}\label{tb:solar}
一列波中$u_E=u_B$, 整个能量密度$u=\eps E^2$. 截面积$A$、长度$c\dd t$的波段在$\dd t$通过截面, 其能量$uAc\dd t$给$S=cu$. 正弦时间平均$\langle E^2\rangle=E_0^2/2$给$I=\eps cE_0^2/2$. 平面波束理想地不随传播距离衰减; 若相同总功率均匀分布到半径$r$球面, 则$I=P/(4\pi r^2)$, 场幅随$1/r$减少. 两者的区别是传播几何而非能量被真空消耗.

\begin{example}
\textbf{原13.1: 太阳常数算例.} 采用原题历史给值$I=1.35\times10^3\,\mathrm{W/m^2}$和日地均距$r=1.50\times10^{11}\,\mathrm m$. 将局部光场近似为正弦平面波, 求场幅和太阳平均辐射功率.
\end{example}
\begin{solution}
$E_0=\sqrt{2I/(\eps c)}\simeq1.01\times10^3\,\mathrm{V/m}$, $B_0=E_0/c\simeq3.36\times10^{-6}\,\mathrm T$. 视太阳各方向近似等功率, $P=4\pi r^2I\simeq3.82\times10^{26}\,\mathrm W$. 这是题设复现, 不是把1350作为当前最新测量. 真实阳光并非一条严格单频波, 等效幅值用于能流量级.
\end{solution}

入射在面积$A$上的能量$\dd U=IA\dd t$携带法向动量$\dd U/c$. 全吸收时$F=IA/c$, 压力$I/c$. 正反射的动量变为$-\dd U/c$, 对物体转移$2\dd U/c$, 压力$2I/c$. 这里$\dd U$是入射能量, 不是反射面所吸收的热. 若正入射部分反射比例$\mathcal R$, 其余全部吸收而无透射, 则压力$(1+\mathcal R)I/c$. 存在透射时还要计透射出射动量; 入射角改变还需投影面积及动量法向分量.

太阳帆的理想反射推力为$2PA/(4\pi r^2c)$, 太阳对帆质量$m$引力为$GM_\odot m/r^2$. 两者比例不依赖$r$, 克服该引力所需面积质量比为$A/m>2\pi cGM_\odot/P$. 这给原13.14.1的推导条件: 正对太阳、理想反射、忽略其他力; 并非任何实际薄膜都自动达到该条件.

\originalsection{电路的能量也穿过周围空间}\label{tb:coax-poynting}
原13.12.4的圆柱导体电导率$\sigma_c$, 半径$a$, 长$\ell$, 稳态均匀电流$I$沿$+z$. 欧姆关系给$E_z=I/(\sigma_c\pi a^2)$, 表面$B_\phi=\mu_0I/(2\pi a)$, 故
\[\vct S(a)=-\frac{I^2}{2\pi^2\sigma_ca^3}\uv r,\qquad
P_{\rm in}=-\oint\vct S\cdot\dd\vct A=\frac{I^2\ell}{\sigma_c\pi a^2}=I^2R.\]
稳态场能不增加, 流入能量变成热. 负号是闭面的向外法向与向内能流相反, 不能删去负号后又把向外积分称为正输入.

同轴传输线内外导体半径$a,b$, 导体间电势差$V$, 轴向电流$I$在内导体向$+z$、外导体返回, 理想绝缘区有
\[E_r=\frac{V}{r\ln(b/a)},\quad B_\phi=\frac{\mu_0I}{2\pi r},\quad
S_z=\frac{VI}{2\pi r^2\ln(b/a)}.\]
整个截面传输功率$\int_a^bS_z2\pi r\dd r=VI$. 主要能流在导体之间的电磁场中, 并非全由电子在金属中携带着能量以漂移速度前行. 实际有损导体还有指入金属的能流, 本理想横截面只算沿轴传递部分.

充电电容的原13.12.3已在本章和第22章逐问推导, 其$\vct S$指向间隙内部使$U_E$增加; 反向放电时相应能流可流出. 长螺线管增加$B$时也从侧面接收能量. 三者共用局部守恒$\partial_tu+\nabla\cdot\vct S=-\jj\cdot\ee$, 不能在有导体耗散的体积内擅自删去右端.

\originalsection{导体反射由表面电流自洽产生}\label{tb:reflection}
本章前面的无入射电流面模型已求规定面电流$K\uv y$向真空两侧发出的横波$E_{\rm rad}=-Z_0K/2$, $Z_0=\sqrt{\mu_0/\eps}$. 现在有沿$+x$入射的$E_i\uv y$照在$x=0$导电面上. 面电流不再预先指定, 而由\emph{总切向场}驱动. 用理想电阻面的面电导$G_s\ge0$表示$K=G_sE_{\rm total}(0)$, $G_s$单位S, 得
\[K=G_s(E_i-Z_0K/2),\quad
K=\frac{G_sE_i}{1+G_sZ_0/2}.\]
若以体电阻率$\rho_r$、厚$D$薄层模拟, 仅在层内场近似均匀时$G_s=D/\rho_r$; 好导体的趋肤条件另须核验, $D\ll\lambda$本身不足保证均匀.

左发射波就是反射波, 右发射波与入射波叠加是透射波. 记$g=Z_0G_s/2$, 振幅比
\[r_E=-\frac g{1+g},\quad t_E=\frac1{1+g}.\]
真空两侧阻抗相同, 功率比例$\mathcal R=g^2/(1+g)^2$, $\mathcal T=1/(1+g)^2$, 吸收比例$\mathcal A=2g/(1+g)^2$. 三者和为1, 面上平均焦耳功率$G_s|E_{\rm total,0}|^2/2$正好等于入射功率乘$\mathcal A$. 最大薄面吸收为$1/2$在$g=1$, 这是该单张对称真空电阻面模型的结论, 不是所有吸收材料的上限.

$G_s\to\infty$时$r_E\to-1,t_E\to0$, 表面切向电场趋零, $K=2E_i/Z_0=2B_i/\mu_0$. 反射波向$-x$, 其$B_z=-E_{r,y}/c$, 因而电场反号、磁场在面上与入射同号相加. 外侧有$B_z=2B_i$, 理想半无限导体内部场为零的边界图像由此成立. 不能在有限导电薄面有部分透射时仍认定右侧场恒零.

\originalsection{驻波有局部能流而无净平均输运}\label{tb:standing-wave}
在完美导体边界, 入射与反射场相加可取
\[E_y=2E_0\sin kx\cos\omega t,\quad
B_z=-\frac{2E_0}c\cos kx\sin\omega t.\]
于是$S_x=-E_0^2\sin(2kx)\sin(2\omega t)/(\mu_0c)$, 一周期平均为零, 但在多数位置瞬时不零, 能量在相邻区域往返. 原例13.2把原点放在电场腹, 写
\[E_y=2E_0\cos kx\cos\omega t,\qquad B_z=\frac{2E_0}c\sin kx\sin\omega t.\]
相应$S_x=+E_0^2\sin(2kx)\sin(2\omega t)/(\mu_0c)$.
两组差别是原点平移, 不是物理矛盾.

两导体面相距$D$形成一维腔, 两端切向电场为节点, $\sin kD=0$, 所以$D=n\lambda/2$, $f_n=nc/(2D)$, $n=1,2,\ldots$. 電磁驻波的电场与磁场在位置上相差$\lambda/4$, 时间上相差$\pi/2$, 不能处处以$B=E/c$把两者当同相行波.

面电流在$t=0$开始变化时, 扰动仅进入$|x|<ct$区域, 波前外尚未收到该变化; 推迟时间$t-|x|/c$直接表达这个因果范围. 有限天线的近区有周期储能往返, 远区才表现为持续外流的辐射. 对短偶极远场$E$近似随$1/r$、强度随$1/r^2$, 方位因子$\sin^2\theta$. 本册用此说明源、波与能流, 不重复高阶电动力学册的完整辐射多极章节.
